\documentclass[10pt,a4paper]{amsart}
\usepackage[margin=19mm]{geometry}
\usepackage{amsmath,amssymb,mathtools}
\usepackage[scr=rsfs,cal=euler]{mathalfa}
\usepackage{xfrac}
\usepackage[unicode,hidelinks,bookmarks=false]{hyperref}
\newcommand{\ie}{i.e.\ }
\newcommand{\cf}{cf.\ }
\newcommand{\resp}{respectively\ }
\newcommand{\Subsection}{\subsection}
\providecommand{\nobreakdash}{\nobreak\hbox{-}}
\setlength{\emergencystretch}{2em}
\multlinegap0pt
\usepackage[shortlabels]{enumitem}
\usepackage{amssymb}
\usepackage{mathtools}
\usepackage{float}
\newtheorem{proposition}{Proposition}
\newtheorem{lemma}{Lemma}
\usepackage{cleveref}
\crefname{proposition}{Proposition}{Propositions}
\crefname{lemma}{Lemma}{Lemmas}
\newcommand{\C}{\mathcal C}
\newcommand{\F}{\mathcal F}
\newcommand{\Hh}{\mathcal H}
\newcommand{\K}{\mathcal K}
\newcommand{\supp}{\operatorname{supp}}
\title{Sunflower-Free Uniform Families: Recursive Constructions and Explicit Bounds}
\author{Edward Axante, Cristian Budala, David Chitic, Bogdan Dumitru, Mihai Nacu}
\date{Source: arXiv:2609.06175v1 (CC BY 4.0)}
\begin{document}
\maketitle

\noindent\emph{Source and licence.} Sunflower-Free Uniform Families: Recursive Constructions and Explicit Bounds. Version arXiv:2609.06175v1 (CC BY 4.0), distributed by arXiv under the Creative Commons Attribution 4.0 International licence, http://creativecommons.org/licenses/by/4.0/, as recorded in the arXiv licence metadata for this exact version and shown on the abstract page https://arxiv.org/abs/2609.06175v1. Full licence notice: \texttt{licenses/arxiv-2609.06175-CC-BY-4.0.txt}.\par\smallskip

\noindent\textit{Editorial scope.} Counted unit: the article's lemma lem:f43-last-profile (source0.tex lines 2027--2030). The article's complete original proof of it is delivered with it (source0.tex lines 2032--2116, one \texttt{proof} environment of 85 lines). No mathematics is written, repaired or re-derived: the statement and the proof are the article's own lines, in the article's own notation, and no retained line is substituted or re-flowed.  Retained uncounted as the source-local context the proof needs: lines 249--253; lines 1821--1833; lines 2018--2021.  Nothing was omitted from any retained span, so the package carries no omission record.  No retained line contains a citation, so the wrapper carries no bibliography and no bibliography entry.  No retained line contains a figure environment, an image-inclusion command or drawing code, so no image is delivered and the package declares no asset dependency.  Editorial formatting only, in addition: an AMS \texttt{amsart} preamble that restates the article's own declarations and macros used by the retained lines, namely the environments \texttt{proposition}, \texttt{lemma} on separate counters, together with the standard packages those lines need.  The retained environments print in this document as Proposition 1, Lemma 1 in order of appearance, whereas the article prints the same items with the numbering of its own full document.  The complete native source of the article as distributed in the arXiv e-print for this exact version is bundled unchanged as \texttt{sources/209499/source0.tex}.
\begin{equation}
 \psi(w,k,m)=\max\bigl\{|\F|:\F\text{ is $w$-uniform, contains no
 $k$-sunflower, and }\nu(\F)\le m-1\bigr\}.
 \label{eq:psi-definition}
\end{equation}

\begin{proposition}[Finite triple bounds]
\label{prop:f43-triple-inputs}
The following statements hold.
\begin{enumerate}[label=\textup{(\roman*)}]
\item Every intersecting, $3$-uniform family without a $3$-sunflower has at
      most $10$ members. In the notation of \cref{eq:psi-definition},
      \[
       \psi(3,3,2)\le10.
      \]
\item If such a family has exactly $10$ members, then its support has at most
      $6$ points.
\end{enumerate}
\end{proposition}

\begin{equation}
 (s,p,t)=(40,12,4).
 \label{eq:f43-last-profile}
\end{equation}

\begin{lemma}[Exclusion of the equality profile]
\label{lem:f43-last-profile}
The profile in \cref{eq:f43-last-profile} cannot occur.
\end{lemma}

\begin{proof}
Assume \cref{eq:f43-last-profile}. Every singleton class is saturated, and
every triple trace is occupied once. Write
\[
 T_m=(R\setminus\{m\})\cup\{z_m\},
 \qquad z_m\notin R
 \quad(m\in R).
\]
For $x\in R$, define
\[
 \K_x=\{F\setminus\{x\}:x\in F\in\F,\ |F\cap R|\ge2\}.
\]

Every member $K\in\K_x$ contains at least one point of
$R\setminus\{x\}$. The support $U_x$ is disjoint from $R$. Hence
$|K\cap U_x|\le2$. If the intersection is a pair, the two design blocks
through that pair, together with $K$, form a sunflower. If it is a singleton
$\{u\}$, the five design blocks through $u$ induce a $2$-regular simple graph
on the other five support points. This graph is a $5$-cycle. Two disjoint
edges of the cycle give two design blocks meeting exactly in $u$; together
with $K$ they form a sunflower. Thus
\begin{equation}
 K\cap U_x=\varnothing\qquad(K\in\K_x).
 \label{eq:f43-K-avoids}
\end{equation}

The family $\K_x$ is intersecting. Two disjoint members, together with any
block of $\Hh_x$, would be three disjoint triples in the link at $x$.
Counting over all $x\in R$ gives
\[
 \sum_{x\in R}|\K_x|=4+2p+3t=40.
\]
Each $\K_x$ is an intersecting triple family without a $3$-sunflower, so
\cref{prop:f43-triple-inputs} bounds it by $10$. Hence every $\K_x$ has
size $10$. By \cref{prop:f43-triple-inputs}(ii), its support has at most six
points. No pair occurs in three members of $\K_x$, since those members would
form a sunflower with that pair as core. The ten triples contain $30$ pair
occurrences. Therefore
\[
 30\le2\binom{|\supp(\K_x)|}{2}\le30.
\]
Equality holds in both bounds. Thus the support has exactly six points, every
pair of support points has degree $2$, and $\K_x$ is a simple
$2$-$(6,3,2)$ design.

Fix $x$ and put $Y=R\setminus\{x\}$. The design $\K_x$ contains
\[
 Y,
 \qquad
 B_m=(Y\setminus\{m\})\cup\{z_m\}
 \quad(m\in Y).
\]
Its support has six points, so the three $z_m$, $m\in Y$, lie in the
three-point complement of $Y$. They are distinct. Indeed, suppose that
$z_i=z_j=q$ and let $y$ be the third point of $Y$. The pairs
$\{y,i\}$, $\{y,j\}$, and $\{y,q\}$ already occur twice among the displayed
blocks. The two remaining design blocks through $y$ must avoid $i,j,q$.
Both would therefore consist of $y$ and the other two outside support points,
contradicting distinctness.

Write $Y=\{y,a,b\}$. The two design blocks through $y$ not displayed above
come from $\C_{\{x,y\}}(R)$. The pair $\{y,z_y\}$ occurs in none of the
displayed blocks, so both new blocks contain $z_y$. The pairs
$\{y,z_a\}$ and $\{y,z_b\}$ each occur once and each needs one further
occurrence. Their residue edges are therefore
\begin{equation}
 \bigl\{\{z_y,z_a\},\{z_y,z_b\}\bigr\}.
 \label{eq:f43-x-view}
\end{equation}
For any two indices in $R$, choose $x$ outside them. The preceding argument
shows that their $z$-points are distinct. Thus all four points $z_m$ are
distinct.

Finally fix distinct $x,y$ and write $R\setminus\{x,y\}=\{a,b\}$. The
$x$-view of the exact pair class is \cref{eq:f43-x-view}. Its $y$-view is
\[
 \bigl\{\{z_x,z_a\},\{z_x,z_b\}\bigr\}.
\]
Both views describe
\[
 \{F\setminus\{x,y\}:F\in\C_{\{x,y\}}(R)\},
\]
so they must be equal. They cannot be equal because the four $z$-points are
distinct.
\end{proof}

\end{document}
